\chapter{Changes to Your Appetite}
\index{Changes to Your Appetite}

\section*{Definition and Overview}
Appetite refers to the desire to eat food, and changes to appetite mean eating noticeably more or less than usual, or losing interest in food altogether. Appetite is regulated by a complex network involving the brain, gut hormones, blood glucose, and signals from the digestive tract, so a wide range of physical and psychological conditions can disturb it. A temporary reduction in appetite during an acute illness is normal and expected, but a persistent change -- whether increased hunger or poor appetite -- may indicate an underlying medical or mental health problem. This chapter outlines the common causes of appetite change, the symptoms that should prompt review, and how appetite problems are investigated and managed.

\section*{Epidemiology}
Changes in appetite are extremely common and can affect people of any age. Reduced appetite is particularly frequent in older adults, in people living with chronic illness, cancer, depression, or dementia, and during recovery from surgery or infection. Increased appetite (hyperphagia) is seen in certain hormonal and metabolic disorders, with some medicines, and in some mental health conditions. Poor appetite in children is a common source of parental worry but is usually temporary and benign. Among older people in hospital and residential care, loss of appetite is a major contributor to malnutrition, frailty, and poorer outcomes. Because appetite is so sensitive to both physical and emotional state, virtually everyone will experience an appetite change at some point in their life.

\section*{Aetiology and Pathophysiology}
\begin{itemize}
    \item \textbf{Infection and acute illness:} fever and inflammation temporarily suppress appetite through effects on the brain's appetite centres
    \item \textbf{Digestive disorders:} nausea, reflux, gastritis, peptic ulcer, coeliac disease, and inflammatory bowel disease make eating uncomfortable
    \item \textbf{Cancer:} tumours can release substances that reduce appetite (anorexia of malignancy), and treatments such as chemotherapy and radiotherapy add to this
    \item \textbf{Endocrine causes:} an underactive thyroid reduces appetite, while overactive thyroid, poorly controlled diabetes, and low blood sugar can increase it
    \item \textbf{Mental health:} depression, anxiety, and stress commonly reduce appetite, whereas some people overeat in response to emotional distress
    \item \textbf{Eating disorders:} anorexia nervosa and bulimia nervosa involve restrictive or disturbed patterns of eating
    \item \textbf{Medicines:} many drugs cause nausea or reduced appetite; others, such as corticosteroids and some antidepressants, stimulate hunger
    \item \textbf{Pregnancy:} hormonal changes can cause nausea, cravings, and food aversions
    \item \textbf{Ageing:} declining sense of taste and smell, slower digestion, and reduced energy needs often lessen appetite
    \item \textbf{Neurological disease:} dementia, stroke, and Parkinson's disease can impair interest in food and the mechanics of eating
\end{itemize}

\section*{Clinical Features}
\begin{itemize}
    \item Eating noticeably less or skipping meals entirely
    \item Loss of interest in foods that were previously enjoyed
    \item Unintentional weight loss or failure to gain weight
    \item Feeling full after only a small amount of food (early satiety)
    \item Constant hunger or eating much more than usual
    \item Strong new cravings for, or aversions to, particular foods
    \item Nausea, bloating, or abdominal discomfort associated with eating
    \item Fatigue, low mood, or disturbed sleep alongside the appetite change
    \item In children, slow growth, irritability, or pickiness around mealtimes
\end{itemize}

\section*{Differential Diagnosis}
\begin{itemize}
    \item Depression, anxiety, and chronic stress
    \item Gastrointestinal conditions, including gastritis, peptic ulcer, coeliac disease, and inflammatory bowel disease
    \item Thyroid disease (hypothyroidism or hyperthyroidism)
    \item Diabetes mellitus, which may cause either reduced appetite or excessive hunger
    \item Cancer, especially of the stomach, pancreas, or lung
    \item Chronic infections such as tuberculosis
    \item Chronic kidney disease or liver disease
    \item Dementia and other neurological disorders
    \item Eating disorders such as anorexia nervosa
    \item Side effects of medicines or substance misuse
    \item Pregnancy-related nausea and cravings
    \item Social factors such as isolation, poverty, or difficulty accessing food
\end{itemize}

\section*{When to Seek Medical Care}
See a doctor if poor appetite lasts more than a week or two, is accompanied by unintentional weight loss, pain, fever, vomiting, or difficulty swallowing, or if it is affecting a child's growth and energy levels. A sudden increase in appetite together with weight loss, excessive thirst, or frequent urination should also be assessed, as these features can indicate diabetes or thyroid disease.

\medskip
\textbf{Contact your local emergency services for:}
\begin{itemize}
    \item Signs of very low blood sugar, such as confusion, sweating, and shaking that do not improve with sweet food or drink
    \item Severe dehydration, with reduced urine output, dizziness, or profound weakness
    \item Sudden inability to swallow, particularly after a stroke-like event
    \item Any sudden, severe, or alarming symptom that worries you
\end{itemize}

\section*{Diagnosis and Investigations}
\begin{itemize}
    \item \textbf{History:} duration of the change, degree of weight change, associated symptoms, current medicines, alcohol use, and mood
    \item \textbf{Physical examination:} weight, height, body mass index, and targeted examination of the abdomen and other systems
    \item \textbf{Blood tests:} full blood count, thyroid function, blood glucose, liver and kidney function, and inflammatory markers
    \item \textbf{Nutritional assessment:} dietary review and testing for vitamin or mineral deficiencies where indicated
    \item \textbf{Imaging and endoscopy:} abdominal ultrasound, gastroscopy, or colonoscopy if a digestive cause is suspected
    \item \textbf{Mental health assessment:} screening for depression, anxiety, and eating disorders
    \item \textbf{Referral:} to a gastroenterologist, endocrinologist, dietitian, or psychologist as appropriate
\end{itemize}

\section*{Management}
\begin{itemize}
    \item \textbf{Treat the underlying cause:} antibiotics for infection, thyroid replacement or suppression, blood-glucose control, antidepressants, or treatment of reflux and nausea
    \item \textbf{Nutritional support:} small, frequent meals, nutrient-dense foods, and dietitian review for people at risk of malnutrition
    \item \textbf{Symptom control:} anti-nausea medicines and treatment of constipation, pain, or mouth problems that interfere with eating
    \item \textbf{Medication review:} checking whether any current medicines are reducing appetite or causing nausea
    \item \textbf{Psychological support:} counselling or cognitive behavioural therapy for depression, anxiety, or eating disorders
    \item \textbf{Oral nutritional supplements:} used under supervision when food intake alone cannot meet requirements
    \item \textbf{Feeding assistance:} practical help with shopping, cooking, and feeding for older or disabled people; discussion of artificial nutrition in advanced illness if appropriate
    \item \textbf{Follow-up:} regular weighing and review to monitor progress and adjust treatment
\end{itemize}

\section*{Complications}
\begin{itemize}
    \item Unintentional weight loss and malnutrition
    \item Vitamin and mineral deficiencies
    \item Muscle wasting, weakness, and increased risk of falls
    \item Impaired immune function and slower recovery from illness
    \item Delayed growth and development in children
    \item Poor wound healing and increased risk of infection in older adults
    \item Worsening of the underlying condition if left untreated
    \item Weight gain, obesity, and metabolic problems when appetite is increased
    \item Social isolation and reduced quality of life
\end{itemize}

\section*{Prognosis and Outcomes}
In most cases, appetite returns to normal once the underlying cause is treated. The temporary appetite loss that accompanies a viral illness typically resolves within days to a week as the person recovers. Where appetite loss results from depression, chronic illness, or cancer, treating the underlying condition and providing nutritional support usually improves intake, although the trajectory depends on the cause. Persistent poor appetite in older or frail people carries a real risk of malnutrition and requires active monitoring. Increased appetite due to hormonal or metabolic causes generally improves when the underlying problem is treated. The outlook in each case should be discussed with the treating clinician.

\section*{Prevention and Self-Care}
\begin{itemize}
    \item Eat regular, balanced meals and keep nutritious foods that you enjoy readily available
    \item Choose small, frequent meals if a large meal feels overwhelming
    \item Drink fluids between meals rather than filling up on them during meals
    \item Manage stress with regular exercise, relaxation, and adequate sleep
    \item Limit alcohol, which can dull the appetite and replace nutritious food
    \item Maintain good dental and mouth care, since pain and poor teeth reduce eating
    \item Monitor weight regularly if appetite is changing
    \item For children, keep mealtimes calm and relaxed and avoid pressuring them to eat
    \item Seek help early for low mood, anxiety, or concerns about eating
    \item Ask a doctor or pharmacist to review medicines if an appetite change follows a new treatment
\end{itemize}

\section*{Sources and Further Reading}
The following reputable sources provide further information on changes in appetite:
\begin{itemize}
    \item NHS. \textit{Health A--Z: symptoms, causes, and self-care advice.} https://www.nhs.uk/conditions/
    \item Mayo Clinic. \textit{Diseases and conditions library.} https://www.mayoclinic.org/
    \item MedlinePlus. \textit{Health topics on nutrition and weight management.} https://medlineplus.gov/
    \item MSD Manual. \textit{Professional and consumer editions on eating and nutrition.} https://www.msdmanuals.com/
    \item World Health Organization. \textit{Nutrition guidance and factsheets.} https://www.who.int/
\end{itemize}
